\section{Aprendizaje por refuerzo condicionado a metas}

\subsection{La meta como descripción de la tarea}

Goal-conditioned RL es el caso particular más importante y más natural del RL multitarea: el identificador de la tarea deja de ser un task ID abstracto y pasa a ser un estado meta $g$. En este caso, la política se escribe $\pi(a\mid s,g)$ y la recompensa se expresa como recompensa por alcanzar la meta o como distancia a la meta:

\begin{importantbox}{Formalización del RL condicionado a metas}
\begin{itemize}
    \item Política: $\pi(a \mid s, g)$
    \item Recompensa dispersa: $r(s,a,g)=1$ si y solo si $d(s,g)<\epsilon$
    \item Recompensa densa: $r(s,a,g)=-d(s,g)$
    \item Objetivo: aprender una política general que vaya desde cualquier punto de partida hacia cualquier meta
\end{itemize}
\end{importantbox}

Una vez que la meta también se considera un task descriptor, muchas de las técnicas del RL multitarea pueden trasladarse directamente a las tareas de goal-reaching.

\subsection{¿Puede reutilizarse la experiencia entre metas?}

Lo verdaderamente interesante de esta sección es que Chelsea Finn no presenta HER como una técnica aislada, sino que primero plantea una pregunta más general. Supongamos que reunimos una experiencia en la tarea 1: ¿puede esa experiencia ayudar también en la tarea 2? Si es así, conviene intentar reinterpretar esa trayectoria con una nueva etiqueta de tarea y una nueva recompensa.

\begin{figure}[H]
\centering
\includegraphics[width=0.9\textwidth]{slides-images/slide-23.jpg}
\caption{La pregunta central del RL multitarea: si una misma trayectoria puede reutilizarse en tareas distintas}
\end{figure}

\begin{figure}[H]
\centering
\includegraphics[width=0.9\textwidth]{slides-images/slide-24.jpg}
\caption{Ejemplo intuitivo: se quería tirar a gol, pero se dio por accidente un pase excelente}
\end{figure}

\subsection{El procedimiento general de Hindsight Relabeling}

\begin{importantbox}{HER (Hindsight Experience Replay)}
La idea central de HER es: \textbf{aunque una trayectoria no cumpla la meta original, puede reetiquetarse con otra meta que sí alcanzó en la práctica.}
\end{importantbox}

El curso divide este proceso en cuatro pasos:
\begin{enumerate}
    \item Recolectar trayectorias con la política actual, $\tau=(s_{1:T}, a_{1:T}, z, r_{1:T})$.
    \item Guardar primero los datos en el replay buffer con la etiqueta de la tarea original.
    \item Elegir después una nueva tarea o meta, recalcular la recompensa y formar una relabeled trajectory.
    \item Usar juntas la trayectoria original y la trayectoria relabeled en la actualización off-policy.
\end{enumerate}

\begin{figure}[H]
\centering
\includegraphics[width=0.9\textwidth]{slides-images/slide-25.jpg}
\caption{Procedimiento unificado de RL multitarea with relabeling: recolectar, reetiquetar y volver a entrenar}
\end{figure}

\begin{knowledgebox}{Cuándo puede hacerse relabeling de forma segura}
En la slide, el curso presenta tres condiciones previas:
\begin{enumerate}
    \item La función de recompensa puede calcularse para la nueva tarea.
    \item La dinámica es la misma, o suficientemente parecida, entre las distintas metas.
    \item El algoritmo de entrenamiento es off-policy, porque una misma experiencia se reutiliza varias veces.
\end{enumerate}
\end{knowledgebox}

\subsection{Por qué goal-conditioned RL es especialmente adecuado para HER}

En el escenario goal-conditioned, el relabeling se simplifica a su forma más elegante: si la trayectoria terminó en $s_T$, se cambia la meta a $g'=s_T$ y, con base en ella, se recalcula la recompensa de toda la trayectoria. Así, una trayectoria fallida que tenía «recompensa cero en todos los pasos» por no haber alcanzado la meta indicada se convierte de inmediato en una muestra exitosa para la nueva meta.

\begin{figure}[H]
\centering
\includegraphics[width=0.9\textwidth]{slides-images/slide-26.jpg}
\caption{Goal-conditioned RL with HER: el estado final alcanzado se toma directamente como nueva meta}
\end{figure}

\begin{figure}[H]
\centering
\includegraphics[width=0.9\textwidth]{slides-images/slide-27.jpg}
\caption{Ejemplo con brazo robótico: desde la perspectiva hindsight, una trayectoria fallida se convierte en una muestra exitosa}
\end{figure}

\begin{warningbox}{Los supuestos implícitos de HER}
HER no puede aplicarse indiscriminadamente en cualquier tarea. Requiere que la recompensa pueda recalcularse para la nueva meta, que las metas compartan la dinámica y que las muestras reetiquetadas del replay buffer no comprometan la estabilidad del entrenamiento. Si la recompensa es muy semántica o no puede reevaluarse, HER no necesariamente es aplicable.
\end{warningbox}

\subsection{Por qué HER es tan eficaz}

Lo más profundo de HER es que cambia nuestra forma de ver los datos de fracaso. Muchas trayectorias de RL resultan «inútiles» no porque carezcan de información, sino porque les asignamos una etiqueta de meta demasiado estrecha. Mediante la reescritura hindsight, HER convierte «no llegué a donde quería ir» en «pero sí llegué a algún lugar», de modo que cada trayectoria aporta supervisión sobre la estructura de alcanzabilidad.

\subsection{Resumen de la sección}

Goal-conditioned RL lleva la idea del RL multitarea a su forma más natural, y HER muestra además que, en tareas con recompensa dispersa, lo verdaderamente escaso no es la experiencia, sino la manera de \textbf{interpretar la experiencia}.
